\documentclass[11pt]{article}
%\usepackage{ctex} %中文
\usepackage[utf8]{inputenc}
\usepackage[a4paper, left=3cm, right=3cm, top=2.5cm, bottom=2.5cm]{geometry} 

\usepackage{graphicx}
\usepackage{subcaption}
\usepackage{float}
\usepackage{amsmath, amsthm, amssymb}

\usepackage{verbatim}
\usepackage{booktabs}
\usepackage{fancyhdr}  % 导入宏包
\setlength{\headheight}{14pt}
\pagestyle{fancy}      % 启用 fancyhdr 页面样式

\fancyhf{}             % 清空原有的页眉和页脚设置
\fancyhead[L]{Assignment 6}
\fancyhead[C]{}
\fancyhead[R]{Algorithms and Analysis}

\fancyfoot[C]{\thepage}

\renewcommand{\headrulewidth}{0.4pt} % 设置页眉线粗细
\renewcommand{\footrulewidth}{0.4pt}

%代码块格式
\usepackage{listings}
\usepackage{xcolor} %设置颜色
%\usepackage{setspace} % 用于调整行间距

%\usepackage{parskip}  

% 定义代码样式：
\lstset{
    language=Python,
    basicstyle=\ttfamily\small,
    numbers=left,
    numberstyle=\tiny\color{gray},
    stepnumber=1,
    breaklines=true,
    frame=shadowbox,
    commentstyle=\color{green!60!black},
    keywordstyle=\color{blue},
    stringstyle=\color{purple},
    backgroundcolor=\color{white},
    showstringspaces=false,
    showspaces=false
}


\title{\textbf{Assignment 06}}
\author{Name: Kaifeng Tang \qquad ID: 524030910124}
\date{\today}

\begin{document}
\maketitle 
\section*{Problem 1} 
\subsection*{(a)} 
\textbf{Transfer to Network: }Build a bipartite graph $G = (L, R, E)$. 
\begin{itemize}
    \item Source node: $s$; Sink node: $t$. 
    \item $L = \{a_1, a_2, \dots, a_k\}$, where $a_i \text{ represents subset }\mathcal{A}_i $. $R = U = \{1, 2, \dots, n\}$. 
    \item Capacity Constraint: $c(e) = 1$ for $\forall e \in E$. 
\end{itemize}
\textbf{Algorithm: }

Function judgeSDR($U$, $\mathcal{A}$): 
\begin{align*}
    & \text{Build flow network } G = (A, U, E) \text{ by rules mentioned above}. \qquad \qquad \qquad \\
    & \text{Run Fax-Flow Algorithm in }G \\
    & \text{get the max-flow } F \\
    & \text{if } F = k: \\
    & \qquad \text{return } true \\
    & \text{else return } false  
\end{align*}
\textbf{Proof of Correctness: }

We need to prove bi-directional corresponding: 
\[\mathcal{A} \text{ has a system of distinct representatives } \iff G \text{exists max-flow }F \]

(i) Sufficiency: assume $F = k$. Since all edges $(s, a_i)$ have capacity 1, 
and there are $k$ edges starting from $s$.

$\implies$ every edge $(a, a_i)$ holds $f(e) = c(e)$, which means full flow. 

$\implies$ each $a_i$ has a flow to some $u_j$. 

Since $(a_i, u_j)$ has capacity only 1, so different $a_i$ must push flow to distinct $u_j$. With capacity 1 in $(u_j, t)$, all flows can sink to $t$. 

$\therefore$ one $u_j$ at most receive flow $f = 1$ from one $a_i$. All these $u_j$ forms a system of distinct representatives. 

(ii) Necessity: assume $\mathcal{A}$ exists a system of distinct representatives. So there exists distinct $u_j \in U$ s.t. $u_{j_i} \in \mathcal{A}_i$. 
So the edges $(a_i, u_j)$ exists. 

We can create a flow: $s \to a_i \to u_j \to t$, push single flow 1 on the path. 
By the capacity constraint, the flow is legal. The total flow: $k$.

On the other hand, the number of edges from $s$ equals to $k$, indicating the max-flow is also $k$.

By Intgrality Theorem, integral capacity $\to$ integral flow. So we have 
\[\mathcal{A} \text{ has a system of distinct representatives } \iff G \text{exists max-flow }F \]
\textbf{Time Complexity: }
Let $m_A = \sum_{i=1}^{k} |\mathcal{A}_i|$. So $|V| = k + n + 2$; $|E| = m_A + k + n$. 

By Ford-Fulkerson Algorithm: flow added 1 by one iteration $\to$ at most $k$ iteration times.
Every time find a path: $O(|E|)$ time. 

$\implies$ Overall time: $O(k|E|)$. 
If using Dinic Algorithm: $O(|V|^2 |E|)$, HKK Algorithm: $O(|E|\sqrt{|V|})$. 

\subsection*{(b)}
\textbf{Transfer to Network: }Build $G' = (V', E')$: with flow path $s \to \mathcal{A}_i \to \mathcal{B}_j \to t$. 
\begin{itemize}
    \item Source node: $s$; Sink node: $t$. 
    \item $u$ means it is the representatives of both an $\mathcal{A}_i$ and an $\mathcal{B}_j$. 
    \item Capacity Constraint: $c(e) = 1$ for $\forall e \in E$. 
\end{itemize}
\textbf{Algorithm: }

Function judgeCommonSDR($U$, $\mathcal{A}$, $\mathcal{B}$): 
\begin{align*}
    & \text{Build flow network } G' = (V', E') \text{ by rules mentioned above}. \qquad \qquad \qquad \\
    & \text{Run Fax-Flow Algorithm in }G' \\
    & \text{get the max-flow } F \\
    & \text{if } F = k: \\
    & \qquad \text{return } true \\
    & \text{else return } false  
\end{align*}
\textbf{Proof of Correctness: }

Claim: $\mathcal{A} \text{ and } \mathcal{B} \text{ has a system of distinct representatives } \iff G \text{exists max-flow }F$

(i) Necessity: assume $\mathcal{A}$ and $\mathcal{B}$ exists a common system of distinct representatives $T = \{u_1, u_2, \dots, u_k\}$.
For $u \in T$: $u$ can represent some $a_i \in \mathcal{A}_i$, $b_j \in \mathcal{B}_j$. 

If we push $f = 1$ on a path: $s \to a_i \to u \to b_j \to t$, by the rules of $G'$, all the edges exist. 
By capacity constraint: $f \le c$, consider a node $u \in T$: $u_{in} = u_{out} \le 1$. So the flow is legal. 
Since the number of edges $(s, a_i)$ = $k$, the max-flow is $k$. 

(ii) Sufficiency: assume $F = k$. Since all edges $(s, a_i)$ have capacity 1, 
and there are $k$ edges starting from $s$, so every edge $(a, a_i)$ holds $f(e) = c(e)$, which means full flow. 

For every $\mathcal{A}_i$: exists some $u$: $s \to a \to u$. Then $u$ can be the representative of $\mathcal{A}$. 

For every $\mathcal{B}_j$: The same. So the same element $u$ can be used by at most one unit flow path.

$\therefore$ all the $u_j$ forms a common system of distinct representatives. 

By Intgrality Theorem, integral capacity $\to$ integral flow. So we have 
\[\mathcal{A} \text{ and } \mathcal{B} \text{ has a system of distinct representatives } \iff G \text{exists max-flow }F\]

\textbf{Time Complexity: }
Let $m_A = \sum_{i=1}^{k} |\mathcal{A}_i|$, $m_B = \sum_{i=1}^{k} |\mathcal{B}_i|$. So $|V'| = k + 2n + 2$; $|E'| = m_A + m_B + 2k + n$. 

The same as Problem (a): Overall time: $O(k|E'|)$. 

If using Dinic Algorithm: $O(|V'|^2 |E'|)$.

HKK Algorithm: $O(|E'|\sqrt{|V'|})$. 

\newpage

\section*{Problem 2}
\subsection*{(a)}
By Dinic's Algorithm: we know that each iteration needs $O(|E|)$ time to find a blocking flow and updating $G^f$. 
Thus, to prove Dinic's Algorithm runs in $O(|E|^{3/2})$ time, we just need to prove Dinic's number of iterations is at most $O(|E|^{1/2})$ time. 

\textbf{\textit{Proof.}} Consider the Algorithm just run $\sqrt{|E|}$ iterations. 

By Property of Dinic's Algorithm: $dist_{i+1}(u) > dist_i (u)$ after one iteration, we can get: 
After $\sqrt{|E|}$ rounds, the distance of all s-t paths in $G^f$ are $\ge \sqrt{|E|}$. 

Let $f = $ Current flow, $f^* = $ max-flow. Consider rest of $|f^*|-|f|$ flows. 

These flows can be decomposed into edge-disjoint s-t paths. Each such path has a length of at least
$\sqrt{|E|}$. Since the number of residual edges in $G^f$ is $O(|E|)$. 

$\implies$ The maximum number of paths with non-intersecting edges: $\frac{O(|E|)}{\sqrt{|E|}}=O(\sqrt{|E|})$. \qed

\subsection*{(b)}
The same as Problem (a): we know each iteration needs $O(|E|)$ time, so we just need to prove 
Dinic's number of iterations is at most $O(|V|^{2/3})$ time. 

\textbf{\textit{Proof.}} Consider Dinic's Algorithm runs $2|V|^{2/3}$ rounds. Suppose Algorithm doesn't halts. 
Suppose that after $2|V|^{2/3}$ rounds, the current flow is $f$. 

Definition: $D_i = \{v \in V: dist_{G^f}(s, v) = i\}$ in $G^f$. Since after each iteration, the shortest distance s-t strictly increases. 

Definition: $\ell = dist_{G^f}(s, t)$

So we have $\ell \ge 2|V|^{2/3}$. 

Now prove $|D_i \cup D_{i+1}| \le |V|^{1/3}$: 

Since each node appears in at most two such unions $D_i \cup D_{i+1} $

$\implies$ $\sum_{i=0}^{\ell - 1}|D_i \cup D_{i+1}| \le 2|V|$, by $\ell \ge 2|V|^{2/3}$ $\implies$ there exists some $i$ s.t. $|D_i \cup D_{i+1}| \le \frac{2|V|}{\ell} \le \frac{2|V|}{2|V|^{2/3}} = |V|^{1/3}$. 

Since $dist(v) \le dist(u) + 1$ for $(u,v)$, so it is necessary to use a residual edge from $D_i \to D_{i+1}$. 
The maximum number of possible edges from $D_i \to D_{i+1}$: 
\[|D_i| \cdot |D_{i+1}| \le |D_i \cup D_{i+1}|^2 \le |V|^{2/3}\]

$\implies$ The number of iterations is at most $O(|V|^{2/3})$. \qed

\section*{Problem 3}
\subsection*{(a)}
The maximum matching problem can be written as an IP problem first:
\begin{align*}
    \text{maximize } \quad 
    & \sum_{e \in E} x_e 
    &&  \\
    \text{subject to } \quad
    & \sum_{e:e \in (u,v)} x_e \le 1
    && (\forall v \in V)\\
    & x_e \in \{0, 1\}
    && (\forall e \in E)
\end{align*}

Objective: represents the number of selected edges $e$, so maximizing $\sum_{e \in E} x_e$ is equivalent to finding the maximum matching.

Constraints: $\sum_{e:e \in (u,v)} x_e \le 1 $ means each vertex can be covered by at most one selected edge, which equals to the definition of matching.
\[x_e = \begin{cases}
    1, & \text{ if edge e is chosen in the matching,} \\
    0, & \text{ otherwise}
\end{cases}\]

In problem (a), constraint 2 is relaxed to $x_e \ge 0$. 
Since any 0/1 solution of a matching satisfies the LP constraint in the problem, 
a feasible solution of a matching must be a feasible solution of the LP.

$\therefore$ This LP is a relaxation of maximum matching problem. 

\subsection*{(b)}
For the vertex constraint: $\sum_{e:e \in (u,v)} x_e \le 1$, introduce a dual variable $y_v \ge 0$ for every vertex $v$. 

The dual objective: $\text{min } \sum_{v \in V} y_v$, since the right side is 1 for primal constraint. 

Dual constraints: $y_v + y_u \ge 1, \quad \forall e \in E$, since the coefficient of primal objective is 1. The constraint holds for every edges $(u, v)$.  

\textbf{The Dual LP:}
\begin{align*}
    \text{minimize } \quad 
    & \sum_{v \in V} y_v 
    &&  \\
    \text{subject to } \quad
    & y_u + y_v \ge 1
    && (\forall e \in E)\\
    & y_v \ge 0 
    && (\forall v \in V)
\end{align*}

\textbf{Explanation: }The dual objective $\sum_{v \in V} y_v $ choose as few vertices as possible, meeting the definition of \textit{Minimum Vertex Cover}. 

Constraint: $y_u + y_v \ge 1$ means that for each edge $e = (u, v)$, 
t least one endpoint of $u$ or $v$ must be selected, so that the edge is covered.

In problem \textit{Minimum Vertex Cover}: $y_v = \begin{cases}
    1, & v \text{ is chosen, } \\
    0, & \text{ otherwise}
\end{cases}$
For dual constraint $y_v \ge 0$, it is a relaxation of IP constraint $y_v \in \{0, 1\}$. 

$\therefore$ The dual LP is a relaxation of Problem \textit{Minimum Vertex Cover}. 

\subsection*{(c)}
Let $G = (L, R, E)$. By bipartite, all $e \in E$ connect a vertex in $L$ and a vertex in $R$; no edge inside $L $ or $R$. 
By definition of \textit{incident matrix}, each column corresponds to one edge in $A$. 
$\implies$ each column has at most two "1"s.

Take any $q \times q$ submatrix $B$ of $A$, we want to prove $\text{det}(B) \in (-1, 0, 1)$ by induction. 

\textit{\textbf{Base.} }When $q = 1$, $B $ has only one element. det$(B) \in (0, 1)$, correct. 

\textit{\textbf{Induction Hypothesis.} }Suppose for all $r \times r$ ($r =1, 2, \dots, q-1$) submatrix $B'$, det$(B) \in (-1, 0, 1) $holds. 

\textit{\textbf{Induction.} }There are only 3 possible cases: 
\begin{itemize}
    \item \textit{Case 1. }There exists a column in $B$ with no 0 $ \to $ det$(B) = 0$, correct. 
    \item \textit{Case 2. }A column has exactly one non-zero element, which is exactly 1 
    $ \to $ By expanding the determinant along this column: det$(B) = \pm \text{det}(B')$, where $B'$ is an incident submatrix with size $(q-1)\times (q-1)$. 
    We can easily get det$(B) \in (-1, 0, 1)$, correct.
    \item \textit{Case 3. }Each column in $B$ contains exactly two 1s. $ \to $ These two 1s must appear in the row corresponding to L and the row corresponding to R, respectively.
    \\ The entire row vector satisfies: $\sum_{\text{row} \in L}\text{row} = \sum_{\text{row} \in R} \text{row}$, which indicates row vectors of $B$ are linearly dependent.
    \\ So det$(B) = 0$, correct. 
\end{itemize}
By conclude, $\text{det}(B) \in (-1, 0, 1)$ holds for all submatrix $B \subseteq A$. 

$\implies$ the \textit{incident matrix} of a bipartite graph is totally unimodular. \qed

\subsection*{(d)}
From 3(b), we can get dual LP of maximum matching. By 3(c), turn LP to matrix form: 
\begin{align*}
    \text{minimize } \quad 
    & \sum_{v \in V} y_v 
    &&  \\
    \text{subject to } \quad
    & A^Ty \ge 1
    && (\forall e \in E)\\
    & y \ge 0 
    && (\forall v \in V)
\end{align*}
where incidence matrix A is totally unimodular.
So we can rewrite Primal LP of maximum matching: 
\begin{align*}
    \text{maximize } \quad 
    & \sum_{e \in E} x_e 
    &&  \\
    \text{subject to } \quad
    & Ax \le 1
    && (\forall v \in V)\\
    & x \ge 0 
    && (\forall e \in E)
\end{align*}

1) Since $A$ is totally unimodular, all poles of a polyhedron $\{x: Ax \le 1, x \ge 0\}$ are integer points $\implies$ 
There exists an integer optimal solution, corresponding exactly to a matching.

2) Since $A$ is totally unimodular, then $A^T$ is totally unimodular. By similar analysis, the dual LP also has an integer optimal solution.
$\implies$ exists an optimal solution with a 0/1 integer, corresponding exactly to a vertex cover. 

By 1), 2), and Strong Dual LP Theorem, Primal LP optimum=Dual LP optimum, 

i.e. maximum matching size = minimum vertex cover size. \qed

\subsection*{(e)}
Consider a graph with triangle structure: 

$G = (V, E)$, where $V = \{1, 2, 3\}$, $E = \{(1,2), (2, 3),(3, 1)\}$. 

Apparently, it is not bipartite. 
\begin{itemize}
    \item Maximum matching: 1
    \item Minimum vertex cover: 2
\end{itemize}
$\implies 1 \ne 2$, the theorem does not hold in non-bipartite graphs.

\section*{Problem 4}
I did this assignment for 5 hours. Difficulty Score: 5. I did it by myself. 
\end{document}



